% А1.3 — Многочлены: теорема Безу, подбор корня, деление уголком, схема Горнера
% Старая методичка: стр. 150–151 (старое задание 15)
% Источники: методичка «Деление уголком» (§ 1, 2, 4, подсказки). Стиль — стиль-изложения.md

\begin{tema}{А1.3}
\opr{Следствие из теоремы Безу.} Если $P(a)=0$, то многочлен $P(x)$ делится на $x-a$:
$P(x)=(x-a)\,Q(x)$, где степень $Q$ на единицу меньше. Так понижают степень
уравнения или неравенства.

\opr{Подбор корня.} Если коэффициенты многочлена целые, то его целые корни нужно искать
среди \emph{делителей свободного члена}. Если к тому же старший коэффициент равен 1,
других рациональных корней нет:
\begin{enumerate}
  \item выпишите делители свободного члена — это кандидаты в корни;
  \item проверяйте их по очереди (подстановкой или схемой Горнера), пока не получите 0;
  \item разделите многочлен на $x-a$ уголком или по схеме Горнера.
\end{enumerate}

\begin{primer}
Решите уравнение $x^3-2x^2+5x-24=0$. Решение: делители числа $-24$: $\pm1,\pm2,\pm3,\dots$
Подходит $x=3$: $27-18+15-24=0$. Разделим многочлен на $x-3$:
\end{primer}

\begin{minipage}[t]{0.5\linewidth}
\opr{Деление уголком}\par\smallskip
\small
$\begin{array}{@{}r@{}r@{}r@{}r|l@{}}
x^3 & {}-2x^2 & {}+5x & {}-24 & \;x-3\\ \cline{5-5}
-(x^3 & {}-3x^2) & & & \;x^2+x+8\\ \cline{1-2}
 & x^2 & {}+5x & & \\
 & -(x^2 & {}-3x) & & \\ \cline{2-3}
 & & 8x & {}-24 & \\
 & & -(8x & {}-24) & \\ \cline{3-4}
 & & & 0 &
\end{array}$
\end{minipage}\hfill
\begin{minipage}[t]{0.46\linewidth}
\opr{Схема Горнера}\par\smallskip
\pfGorner{3}{1,-2,5,-24}\par\smallskip
\footnotesize\raggedright
\begin{enumerate}[leftmargin=4mm, itemsep=0pt, topsep=0pt, label=\arabic*)]
  \item \textcolor{pfBrend}{сносим вниз} $1$;
  \item \textcolor{pfOrange}{$1\cdot3=3$}, \ $-2+3=1$;
  \item \textcolor{pfOrange}{$1\cdot3=3$}, \ $5+3=8$;
  \item \textcolor{pfOrange}{$8\cdot3=24$}, \ $-24+24=0$.
\end{enumerate}
Умножаем по стрелке, складываем по столбцу. Остаток~0, значит, 3~корень; остальные числа~— коэффициенты частного.
\end{minipage}

\medskip
Итог: $x^3-2x^2+5x-24=(x-3)(x^2+x+8)$. У $x^2+x+8$ дискриминант
$D=1-32<0$, поэтому $x=3$ — единственный корень. Ответ: 3.
\end{tema}
